\chapter{Reiterman Theory for Effective Input Quotients of Residual Mealy Recognizers}
\label{ch:reiterman-effective-input-quotients-residual-mealy}

\section{Overview}
\label{sec:reiterman-overview}

The preceding chapter established the state-level Eilenberg theory for
residual Mealy recognition. For internal categories
\[
\mathbb A
\qquad\text{and}\qquad
\mathbb B
\]
in a Grothendieck topos \(\mathcal E\), it studied behaviour specifications
\[
S\hookrightarrow \mathsf{Beh}_{\mathbb A,\mathbb B}
\]
through their residual closures
\[
\mathsf D(S)\hookrightarrow \mathsf{Beh}_{\mathbb A,\mathbb B}.
\]
The object \(\mathsf D(S)\) is the minimal separated residual state recognizer
generated by \(S\). The corresponding residual syntactic state category is
\[
\operatorname{SynState}_{\mathbb A,\mathbb B}(S)
:=
\int_{\mathbb A}\mathsf D(S).
\]

This chapter develops the quotient-level theory. The passage from the
Eilenberg chapter to the Reiterman chapter is the passage
\[
\text{closed residual state recognizer}
\quad\longmapsto\quad
\text{effective input quotient}.
\]
For a closed behaviour system
\[
D\hookrightarrow\mathsf{Beh}_{\mathbb A,\mathbb B},
\]
the Mealy action of input arrows on \(D\) determines, for every arrow of
\(\mathbb A\), a transition-output operation on all residual states of \(D\).
Because Mealy execution processes input arrows target-to-source, the relevant
input category is
\[
F:=\mathbb A^{\mathrm{op}}.
\]
The central representation is an internal functor
\[
\rho_D:
\mathbb A^{\mathrm{op}}
\longrightarrow
\operatorname{TransOut}_{\mathbb B}(D),
\]
where
\[
\operatorname{TransOut}_{\mathbb B}(D)
\]
is the internal category of transition-output endomorphisms of \(D\).

The algebraic recognizer of \(D\) is the image of this functor:
\[
\operatorname{Alg}_{\mathbb A,\mathbb B}(D)
:=
\operatorname{im}(\rho_D).
\]
For a behaviour specification \(S\), the syntactic algebraic recognizer is
\[
\operatorname{SynAlg}_{\mathbb A,\mathbb B}(S)
:=
\operatorname{Alg}_{\mathbb A,\mathbb B}(\mathsf D(S)).
\]
Equivalently,
\[
\operatorname{SynAlg}_{\mathbb A,\mathbb B}(S)
=
\operatorname{im}
\left(
\mathbb A^{\mathrm{op}}
\to
\operatorname{TransOut}_{\mathbb B}(\mathsf D(S))
\right).
\]

Thus the full Eilenberg--Reiterman picture has three levels:
\[
\text{behaviour specifications}
\quad\longleftrightarrow\quad
\text{residual state recognizers}
\quad\longrightarrow\quad
\text{effective input quotients}.
\]
The Eilenberg correspondence belongs to the first two levels. The Reiterman
correspondence belongs to the third.

The quotient-level theory is first local in the fixed input category
\[
F=\mathbb A^{\mathrm{op}}.
\]
Local effective input quotients are identity-on-objects regular quotients
\[
F\twoheadrightarrow C
\]
of internal categories, or equivalently effective internal category
congruences on \(F\).

There are two Reiterman compactness regimes.

The first is an abstract intersection-compact regime. For a regular cardinal
\(\kappa\), quotient classes of \(\kappa\)-intersection-compact input
quotients closed under further compact quotients and \(<\kappa\)-subdirect
products are exactly the classes definable by completed input equations in
the corresponding compact completion.

The second is the finite/profinite regime. In ordinary Set-based finite
recognition, finite quotients are controlled not by meet-compactness of their
kernel congruences, but by compactness of profinite completions and finite
determinacy of continuous maps into finite discrete quotients. Accordingly,
this chapter treats finite/profinite Reiterman theory as a separate theorem,
not as a special case of intersection compactness.

Finally, the chapter globalizes the quotient theory across varying Mealy
interfaces. Reiterman-admissible interface translations are required to
preserve residual closure, algebraic recognizers, quotient reindexing, and
completed equations. Under these hypotheses, local Reiterman theorems assemble
into global quotient-equational correspondences.

\section{Quotient-level ambient setting}
\label{sec:reiterman-ambient-setting}

Throughout the internal-categorical quotient sections of this chapter, assume
that
\[
\mathcal E
\]
is a Grothendieck topos. Thus \(\mathcal E\) is locally cartesian closed,
regular, Barr-exact, and locally presentable. Every slice
\[
\mathcal E/Z
\]
is again regular and Barr-exact.

Let
\[
\mathbb A
=
(A_0,A_1,s_A,t_A,e_A,m_A)
\qquad\text{and}\qquad
\mathbb B
=
(B_0,B_1,s_B,t_B,e_B,m_B)
\]
be internal categories in \(\mathcal E\). We continue to use the composition
convention from the preceding chapters: if
\[
a:u\to v,
\qquad
b:v\to w,
\]
then
\[
b\circ a:u\to w
\]
means first \(a\), then \(b\). Equivalently,
\[
m_A(a,b)=b\circ a.
\]

Write
\[
Z:=A_0\times B_0
\]
and
\[
\mathsf{Beh}
:=
\mathsf{Beh}_{\mathbb A,\mathbb B}
\]
for the final behaviour system for Mealy morphisms
\[
\mathbb A\rightsquigarrow\mathbb B.
\]

Throughout this chapter,
\[
D\hookrightarrow\mathsf{Beh}
\]
denotes a closed behaviour system, equivalently a separated residual state
recognizer. Its structure maps are
\[
p_D:D\to A_0,
\qquad
q_D:D\to B_0.
\]
Its transition and output maps are the restrictions of those of
\(\mathsf{Beh}\):
\[
\delta_D:
A_1\times_{t_A,A_0,p_D}D\to D,
\]
and
\[
\omega_D:
A_1\times_{t_A,A_0,p_D}D\to B_1.
\]

Because Mealy execution processes an input arrow
\[
a:u\to v
\]
at a state over \(v\) and moves to a state over \(u\), the input category for
quotient theory is
\[
F:=\mathbb A^{\mathrm{op}}.
\]
Thus \(F\) has object object \(A_0\), and an arrow
\[
a:u\to v
\]
of \(\mathbb A\) is regarded as an arrow
\[
a^{\mathrm{op}}:v\to u
\]
of \(F\).

The goal is to associate to a closed residual recognizer \(D\) an effective
input quotient
\[
F\twoheadrightarrow \operatorname{Alg}_{\mathbb A,\mathbb B}(D)
\]
which identifies precisely those input arrows that induce the same
transition-output operation on all states of \(D\).

\section{The transition-output endomorphism category}
\label{sec:reiterman-transout-category}

Let
\[
D\hookrightarrow\mathsf{Beh}
\]
be a closed behaviour system. An input arrow of \(\mathbb A\) acts on the
fibres of \(D\) and emits an output arrow of \(\mathbb B\). We first construct
the internal category whose arrows are all possible such transition-output
operations.

Informally, an arrow of
\[
\operatorname{TransOut}_{\mathbb B}(D)
\]
from
\[
v\in A_0
\]
to
\[
u\in A_0
\]
is a pair
\[
(\tau,o)
\]
where
\[
\tau:D_v\to D_u
\]
is a transition function on residual states, and where, for each state
\[
x\in D_v,
\]
there is an output arrow
\[
o_x:q_D(\tau x)\to q_D(x)
\]
in \(\mathbb B\).

We now construct this category internally.

Let
\[
R:=A_0\times A_0.
\]
We write a generalized point of \(R\) as \((u,v)\), where \(v\) is the source
input object and \(u\) is the target input object. Let
\[
\pi_t,\pi_s:R\to A_0
\]
denote the first and second projections:
\[
\pi_t(u,v)=u,
\qquad
\pi_s(u,v)=v.
\]
Define objects of \(\mathcal E/R\)
\[
D_{\mathrm{src}}
:=
R\times_{\pi_s,A_0,p_D}D,
\qquad
D_{\mathrm{tgt}}
:=
R\times_{\pi_t,A_0,p_D}D.
\]
Thus over \((u,v)\), these are respectively
\[
D_v
\qquad\text{and}\qquad
D_u.
\]

Because \(\mathcal E\) is locally cartesian closed, the slice
\[
\mathcal E/R
\]
is cartesian closed. Let
\[
H:=[D_{\mathrm{src}},D_{\mathrm{tgt}}]_{\mathcal E/R}
\]
be the internal hom in \(\mathcal E/R\). It comes with an evaluation map
\[
\operatorname{ev}_{\tau}:
H\times_R D_{\mathrm{src}}
\to
D_{\mathrm{tgt}}.
\]
A generalized point of \(H\) over \((u,v)\) is a map
\[
\tau:D_v\to D_u.
\]

Let
\[
U:=H\times_R D_{\mathrm{src}}.
\]
There are two maps
\[
U\to B_0.
\]
The first is
\[
q_{\mathrm{out}}
:=
q_D\operatorname{ev}_{\tau}:U\to B_0,
\]
and the second is
\[
q_{\mathrm{in}}
:=
q_D\pi_{D_{\mathrm{src}}}:U\to B_0.
\]
Define
\[
P:=
U\times_{(q_{\mathrm{out}},q_{\mathrm{in}}),B_0\times B_0,(s_B,t_B)}B_1.
\]
Thus a generalized point of \(P\) over \((\tau,x)\) is an arrow
\[
q_D(\tau x)\to q_D(x)
\]
of \(\mathbb B\).

Let
\[
\pi_U:U\to H
\]
be the projection. Since \(\mathcal E/H\) is locally cartesian closed, the
object of sections of \(P\to U\) over \(H\) exists. We denote it by
\[
\operatorname{Sec}_H(P)\to H.
\]
Equivalently,
\[
\operatorname{Sec}_H(P)
\]
is the object over \(H\) representing maps
\[
o:U\to B_1
\]
satisfying
\[
s_Bo=q_{\mathrm{out}},
\qquad
t_Bo=q_{\mathrm{in}}.
\]

\begin{definition}[Transition-output endomorphism category]
\label{def:reiterman-transout-category}
The \textbf{transition-output endomorphism category} of \(D\), denoted
\[
\operatorname{TransOut}_{\mathbb B}(D),
\]
is the internal category defined as follows.

Its object object is
\[
A_0.
\]

Its arrow object is
\[
\operatorname{TransOut}_{\mathbb B}(D)_1
:=
\operatorname{Sec}_H(P),
\]
viewed over
\[
R=A_0\times A_0.
\]
The source and target maps are the composites with
\[
\pi_s,\pi_t:R\to A_0.
\]
Thus a generalized arrow
\[
(\tau,o):v\to u
\]
consists of a map
\[
\tau:D_v\to D_u
\]
and an internally parametrized output assignment
\[
o_x:q_D(\tau x)\to q_D(x)
\]
for generalized states \(x\in D_v\).

The identity at \(v\in A_0\) is
\[
(\operatorname{id}_{D_v}, e_Bq_D(-)).
\]

Given arrows
\[
(\tau,o):v\to u
\qquad\text{and}\qquad
(\sigma,r):u\to w,
\]
their composite
\[
(\sigma,r)\circ(\tau,o):v\to w
\]
has transition part
\[
\sigma\tau:D_v\to D_w
\]
and output component at \(x\in D_v\)
\[
q_D(\sigma\tau x)
\xrightarrow{r_{\tau x}}
q_D(\tau x)
\xrightarrow{o_x}
q_D(x).
\]
Equivalently,
\[
((\sigma,r)\circ(\tau,o))_x
=
o_x\circ r_{\tau x}.
\]
\end{definition}

\begin{proposition}
\label{prop:reiterman-transout-is-internal-category}
The data of Definition~\ref{def:reiterman-transout-category} define an
internal category
\[
\operatorname{TransOut}_{\mathbb B}(D)
\]
with object object \(A_0\).
\end{proposition}

\begin{proof}
The identity map is induced by the identity transition
\[
D_v\to D_v
\]
and the identity output section
\[
x\longmapsto e_B(q_Dx).
\]

The composition map is induced by composition in the internal homs and by
composition in \(\mathbb B\). For composable generalized arrows
\[
(\tau,o):v\to u
\qquad\text{and}\qquad
(\sigma,r):u\to w,
\]
the transition part is
\[
\sigma\tau:D_v\to D_w.
\]
The output part is well typed because
\[
r_{\tau x}:q_D(\sigma\tau x)\to q_D(\tau x)
\]
and
\[
o_x:q_D(\tau x)\to q_D(x),
\]
so their composite is an arrow
\[
q_D(\sigma\tau x)\to q_D(x)
\]
of \(\mathbb B\).

The identity laws follow from the identity laws in \(\mathbb B\) and for maps
between fibres of \(D\). Associativity of the transition parts follows from
associativity of composition of maps. Associativity of the output parts
follows from associativity of composition in \(\mathbb B\):
\[
o_x\circ(r_{\tau x}\circ s_{\sigma\tau x})
=
(o_x\circ r_{\tau x})\circ s_{\sigma\tau x}.
\]
Therefore the internal-category laws hold.
\end{proof}

\section{The canonical transition-output representation}
\label{sec:reiterman-canonical-transition-output-representation}

Let
\[
D\hookrightarrow\mathsf{Beh}
\]
be a closed behaviour system. Its Mealy structure gives an input action
\[
\delta_D:
A_1\times_{t_A,A_0,p_D}D\to D
\]
and an output map
\[
\omega_D:
A_1\times_{t_A,A_0,p_D}D\to B_1.
\]

An input arrow
\[
a:u\to v
\]
acts on states over \(v\), sends them to states over \(u\), and produces an
output arrow in \(\mathbb B\). Thus \(a\) determines an arrow
\[
v\to u
\]
in
\[
\operatorname{TransOut}_{\mathbb B}(D).
\]
Therefore the representation is a functor out of
\[
\mathbb A^{\mathrm{op}}.
\]

\begin{definition}[Canonical transition-output representation]
\label{def:reiterman-rho-D}
The \textbf{canonical transition-output representation} of \(D\) is the
internal functor
\[
\rho_D:
\mathbb A^{\mathrm{op}}
\to
\operatorname{TransOut}_{\mathbb B}(D)
\]
defined as follows.

On objects,
\[
\rho_{D,0}=\operatorname{id}_{A_0}.
\]
On an arrow
\[
a:u\to v
\]
of \(\mathbb A\), regarded as an arrow
\[
a^{\mathrm{op}}:v\to u
\]
of \(\mathbb A^{\mathrm{op}}\), define
\[
\rho_D(a)=(\tau_a,o_a),
\]
where
\[
\tau_a(x):=\delta_D(a,x)
\]
and
\[
o_a(x):=\omega_D(a,x).
\]
\end{definition}

\begin{proposition}
\label{prop:reiterman-rho-D-functor}
The map
\[
\rho_D:
\mathbb A^{\mathrm{op}}
\to
\operatorname{TransOut}_{\mathbb B}(D)
\]
is an internal functor.
\end{proposition}

\begin{proof}
For identities, let \(x\in D_v\). The Mealy unit laws give
\[
\delta_D(e_A(v),x)=x
\]
and
\[
\omega_D(e_A(v),x)=e_B(q_Dx).
\]
Thus
\[
\rho_D(e_A(v))
\]
is the identity transition-output operation at \(v\).

Now let
\[
a:u\to v,
\qquad
b:v\to w
\]
be composable arrows in \(\mathbb A\). In
\[
\mathbb A^{\mathrm{op}},
\]
one has
\[
a^{\mathrm{op}}\circ b^{\mathrm{op}}
=
(b\circ a)^{\mathrm{op}}.
\]
Thus functoriality requires
\[
\rho_D(b\circ a)
=
\rho_D(a)\circ\rho_D(b).
\]

Let \(x\in D_w\). The transition part of
\[
\rho_D(b\circ a)
\]
is
\[
\delta_D(b\circ a,x).
\]
By the Mealy transition law,
\[
\delta_D(b\circ a,x)
=
\delta_D(a,\delta_D(b,x)).
\]
This is the transition part of
\[
\rho_D(a)\circ\rho_D(b).
\]

The output part of
\[
\rho_D(b\circ a)
\]
at \(x\) is
\[
\omega_D(b\circ a,x).
\]
By the Mealy output multiplication law,
\[
\omega_D(b\circ a,x)
=
\omega_D(b,x)\circ
\omega_D(a,\delta_D(b,x)).
\]
This is precisely the output composition rule in
\[
\operatorname{TransOut}_{\mathbb B}(D).
\]
Therefore
\[
\rho_D(b\circ a)=\rho_D(a)\circ\rho_D(b),
\]
and \(\rho_D\) is an internal functor.
\end{proof}

\section{Algebraic recognizers and syntactic input congruences}
\label{sec:reiterman-algebraic-recognizers}

We now define the algebraic recognizer canonically associated to a closed
residual state recognizer.

\begin{lemma}[Image of an identity-on-objects internal functor]
\label{lem:reiterman-image-identity-on-objects-functor}
Let
\[
F:\mathcal C\to\mathcal D
\]
be an internal functor in a regular category, and suppose that
\[
F_0:\mathcal C_0\to\mathcal D_0
\]
is the identity on a common object object \(O\). Then the image of
\[
F_1:\mathcal C_1\to\mathcal D_1
\]
carries a unique internal-category structure making
\[
\mathcal C
\twoheadrightarrow
\operatorname{im}(F)
\hookrightarrow
\mathcal D
\]
a factorization by an identity-on-objects regular epimorphism followed by an
identity-on-objects monomorphism.
\end{lemma}

\begin{proof}
Let
\[
\mathcal C_1\twoheadrightarrow I\hookrightarrow\mathcal D_1
\]
be the image factorization of \(F_1\). Since \(F\) preserves source and
target, the subobject \(I\hookrightarrow \mathcal D_1\) lies over
\[
O\times O.
\]
The identities of \(\mathcal D\) land in \(I\) because \(F\) preserves
identities and \(F_0\) is the identity on objects.

It remains to show closure under composition. Pull back the composable-pair
object
\[
I\times_O I
\]
along the regular epimorphism
\[
\mathcal C_1\twoheadrightarrow I.
\]
Regular epimorphisms are stable under pullback, so the induced map from the
composable-pair object of \(\mathcal C\) onto \(I\times_O I\) is a regular
epimorphism. After this pullback, composition in \(I\) is represented by
composition in \(\mathcal C\), followed by \(F_1\), and therefore lands in
\(I\). Since pullback along a regular epimorphism reflects containment of
subobjects, composition in \(\mathcal D\) restricts to a map
\[
I\times_O I\to I.
\]

The internal-category axioms hold because they hold in \(\mathcal D\) and
because
\[
I\hookrightarrow \mathcal D_1
\]
is monic. Uniqueness is forced by the requirement that
\[
I\hookrightarrow\mathcal D_1
\]
be an internal functor.
\end{proof}

\begin{definition}[Algebraic recognizer of a closed behaviour system]
\label{def:reiterman-alg-D}
Let
\[
D\hookrightarrow\mathsf{Beh}
\]
be a closed behaviour system. The \textbf{algebraic recognizer} of \(D\) is
the image of the canonical transition-output representation
\[
\rho_D:
\mathbb A^{\mathrm{op}}
\to
\operatorname{TransOut}_{\mathbb B}(D).
\]
It is denoted
\[
\operatorname{Alg}_{\mathbb A,\mathbb B}(D).
\]
Thus
\[
\rho_D
\]
factors as
\[
\mathbb A^{\mathrm{op}}
\twoheadrightarrow
\operatorname{Alg}_{\mathbb A,\mathbb B}(D)
\hookrightarrow
\operatorname{TransOut}_{\mathbb B}(D),
\]
where the first functor is identity on objects and regular epimorphic on
arrows, and the second is identity on objects and monic on arrows.
\end{definition}

\begin{definition}[Syntactic algebraic recognizer]
\label{def:reiterman-synalg}
For a behaviour specification
\[
S\hookrightarrow\mathsf{Beh},
\]
the \textbf{syntactic algebraic recognizer} of \(S\) is
\[
\operatorname{SynAlg}_{\mathbb A,\mathbb B}(S)
:=
\operatorname{Alg}_{\mathbb A,\mathbb B}(\mathsf D(S)).
\]
Equivalently,
\[
\operatorname{SynAlg}_{\mathbb A,\mathbb B}(S)
=
\operatorname{im}
\left(
\mathbb A^{\mathrm{op}}
\to
\operatorname{TransOut}_{\mathbb B}(\mathsf D(S))
\right).
\]
\end{definition}

\begin{definition}[Syntactic input congruence]
\label{def:reiterman-syntactic-input-congruence}
Let
\[
D\hookrightarrow\mathsf{Beh}
\]
be a closed behaviour system. The \textbf{syntactic input congruence}
\[
{\equiv_D}
\]
on
\[
\mathbb A^{\mathrm{op}}
\]
is the kernel pair of the arrow map of
\[
\rho_D:
\mathbb A^{\mathrm{op}}
\to
\operatorname{TransOut}_{\mathbb B}(D).
\]
In generalized-element notation, for parallel arrows
\[
a,a':u\to v
\]
of \(\mathbb A\), one has
\[
a\equiv_D a'
\]
if and only if, for every state
\[
x\in D_v,
\]
one has
\[
\delta_D(a,x)=\delta_D(a',x)
\]
and
\[
\omega_D(a,x)=\omega_D(a',x).
\]
Equivalently,
\[
a\equiv_D a'
\]
if and only if
\[
\rho_D(a)=\rho_D(a')
\]
as arrows of
\[
\operatorname{TransOut}_{\mathbb B}(D).
\]
\end{definition}

\begin{proposition}
\label{prop:reiterman-syntactic-input-congruence}
The relation
\[
{\equiv_D}
\]
is an internal category congruence on
\[
\mathbb A^{\mathrm{op}}.
\]
Moreover,
\[
\operatorname{Alg}_{\mathbb A,\mathbb B}(D)
\cong
\mathbb A^{\mathrm{op}}/{\equiv_D}.
\]
\end{proposition}

\begin{proof}
The relation
\[
{\equiv_D}
\]
is the kernel pair of the arrow map of the internal functor
\[
\rho_D:
\mathbb A^{\mathrm{op}}
\to
\operatorname{TransOut}_{\mathbb B}(D).
\]
Therefore it is an internal equivalence relation compatible with source and
target. Compatibility with identities and composition follows because it is
the kernel pair of an internal functor.

Since \(\mathcal E\) is exact, the quotient internal category exists. Because
\[
{\equiv_D}
\]
is the kernel pair of \(\rho_D\), the quotient is canonically the image of
\(\rho_D\), namely
\[
\operatorname{Alg}_{\mathbb A,\mathbb B}(D).
\]
\end{proof}

\section{Algebraic action categories}
\label{sec:reiterman-algebraic-action-categories}

The residual state category
\[
\int_{\mathbb A}D
\]
uses the original input category \(\mathbb A\). The algebraic recognizer
\[
\operatorname{Alg}_{\mathbb A,\mathbb B}(D)
\]
acts on \(D\) through its inclusion into
\[
\operatorname{TransOut}_{\mathbb B}(D).
\]
This yields a quotient state/process category in which input arrows are
identified precisely when they induce the same transition-output operation
on \(D\).

\begin{definition}[Action category of the algebraic recognizer]
\label{def:reiterman-algebraic-action-category}
Let
\[
D\hookrightarrow\mathsf{Beh}
\]
be a closed behaviour system. The inclusion
\[
\operatorname{Alg}_{\mathbb A,\mathbb B}(D)
\hookrightarrow
\operatorname{TransOut}_{\mathbb B}(D)
\]
determines an action of
\[
\operatorname{Alg}_{\mathbb A,\mathbb B}(D)
\]
on the state object \(D\).

Explicitly, if
\[
\alpha:v\to u
\]
is an arrow of
\[
\operatorname{Alg}_{\mathbb A,\mathbb B}(D),
\]
its image in
\[
\operatorname{TransOut}_{\mathbb B}(D)
\]
has transition component
\[
\tau_\alpha:D_v\to D_u
\]
and output component
\[
o_\alpha(x):q_D(\tau_\alpha x)\to q_D(x).
\]
The action is
\[
\alpha\cdot x:=\tau_\alpha(x).
\]

The corresponding action category is denoted
\[
\int_{\operatorname{Alg}(D)}D.
\]
Its object object is
\[
D.
\]
Its arrow object consists of pairs
\[
(\alpha,x),
\]
where
\[
\alpha:v\to u
\]
is an arrow of
\[
\operatorname{Alg}_{\mathbb A,\mathbb B}(D)
\]
and
\[
x\in D_v.
\]
Its source and target are
\[
s(\alpha,x)=x,
\qquad
t(\alpha,x)=\alpha\cdot x.
\]
The output labelling sends
\[
(\alpha,x)
\]
to
\[
o_\alpha(x):q_D(\alpha\cdot x)\to q_D(x).
\]
\end{definition}

\begin{proposition}[State category quotients to algebraic action category]
\label{prop:reiterman-state-to-algebraic-quotient}
Let
\[
D\hookrightarrow\mathsf{Beh}
\]
be a closed behaviour system. The quotient
\[
\mathbb A^{\mathrm{op}}
\twoheadrightarrow
\operatorname{Alg}_{\mathbb A,\mathbb B}(D)
\]
induces a labelled regular quotient of action categories
\[
\int_{\mathbb A}D
\twoheadrightarrow
\int_{\operatorname{Alg}(D)}D.
\]
Fibrewise, it sends
\[
(a,x)
\]
to
\[
([a],x),
\]
where \([a]\) denotes the \({\equiv_D}\)-class of \(a\).
\end{proposition}

\begin{proof}
The quotient
\[
\mathbb A^{\mathrm{op}}
\twoheadrightarrow
\operatorname{Alg}_{\mathbb A,\mathbb B}(D)
\]
identifies precisely those input arrows that induce equal transition-output
operations on all states of \(D\). Therefore, for each state
\[
x\in D,
\]
the transition
\[
\delta_D(a,x)
\]
and output
\[
\omega_D(a,x)
\]
depend only on the class \([a]\). This gives a well-defined internal functor
\[
\int_{\mathbb A}D
\to
\int_{\operatorname{Alg}(D)}D.
\]
It is the pullback, along the state projection \(D\to A_0\), of the
identity-on-objects regular epi
\[
\mathbb A^{\mathrm{op}}
\twoheadrightarrow
\operatorname{Alg}_{\mathbb A,\mathbb B}(D)
\]
on arrow objects, and is therefore regular epi on arrows and identity on
objects. Source, target, identity, composition, input labelling, and output
labelling are preserved because the quotient is induced by the functor
\[
\rho_D.
\]
\end{proof}

\section{Algebraic syntactic divisibility}
\label{sec:reiterman-algebraic-syntactic-divisibility}

Let
\[
X:\mathbb A\rightsquigarrow\mathbb B
\]
be a Mealy system and write
\[
D_X:=\operatorname{Real}(X).
\]
By the results of the preceding chapter, \(D_X\) is a closed behaviour system,
and hence has an algebraic recognizer
\[
\operatorname{Alg}_{\mathbb A,\mathbb B}(D_X).
\]

\begin{theorem}[Algebraic syntactic divisibility]
\label{thm:reiterman-algebraic-syntactic-divisibility}
Let
\[
S\hookrightarrow\mathsf{Beh}
\]
be a behaviour specification. If a Mealy system
\[
X:\mathbb A\rightsquigarrow\mathbb B
\]
recognizes \(S\), then
\[
\operatorname{SynAlg}_{\mathbb A,\mathbb B}(S)
\]
is a regular quotient of
\[
\operatorname{Alg}_{\mathbb A,\mathbb B}(\operatorname{Real}(X)).
\]
\end{theorem}

\begin{proof}
Since \(X\) recognizes \(S\), the state Eilenberg theorem gives
\[
\mathsf D(S)\leq\operatorname{Real}(X).
\]
Let
\[
D:=\mathsf D(S),
\qquad
E:=\operatorname{Real}(X).
\]
Then
\[
D\leq E
\]
is an inclusion of closed behaviour systems.

Every input-induced transition-output operation on \(E\) restricts to one on
\(D\), because \(D\) is closed under the input action. Equivalently, let
\[
\theta_E
\qquad\text{and}\qquad
\theta_D
\]
be the input congruences on
\[
\mathbb A^{\mathrm{op}}
\]
induced by \(E\) and \(D\), respectively. If two input arrows induce the same
transition-output operation on all of \(E\), then they induce the same
transition-output operation on the subobject \(D\). Hence
\[
\theta_E\leq\theta_D.
\]
Therefore there is an identity-on-objects regular quotient
\[
\mathbb A^{\mathrm{op}}/\theta_E
\twoheadrightarrow
\mathbb A^{\mathrm{op}}/\theta_D.
\]
Using
\[
\mathbb A^{\mathrm{op}}/\theta_E
\cong
\operatorname{Alg}_{\mathbb A,\mathbb B}(E)
\]
and
\[
\mathbb A^{\mathrm{op}}/\theta_D
\cong
\operatorname{SynAlg}_{\mathbb A,\mathbb B}(S),
\]
the result follows.
\end{proof}

\section{Effective input quotients}
\label{sec:reiterman-effective-input-quotients}

The Reiterman side is formulated in terms of quotients of the fixed input
category
\[
F:=\mathbb A^{\mathrm{op}}.
\]
We work with internal categories whose object object is \(A_0\) and with
identity-on-objects internal functors.

\begin{definition}[Internal category congruence]
\label{def:reiterman-internal-category-congruence}
An \textbf{internal category congruence} on
\[
F
\]
is an internal equivalence relation
\[
\theta\hookrightarrow
F_1\times_{A_0\times A_0}F_1
\]
on the arrow object of \(F\), over the common source and target object
\[
A_0\times A_0,
\]
compatible with identities and composition.

Compatibility with identities says that identity arrows are related only to
identity arrows over the same object. Compatibility with composition says
that whenever
\[
f\sim_\theta f'
\qquad\text{and}\qquad
g\sim_\theta g'
\]
are composable pairs, then
\[
g\circ f\sim_\theta g'\circ f'.
\]
\end{definition}

\begin{definition}[Effective input quotient]
\label{def:reiterman-effective-input-quotient}
An \textbf{effective input quotient} of
\[
F=\mathbb A^{\mathrm{op}}
\]
is an identity-on-objects internal functor
\[
q:F\twoheadrightarrow C
\]
whose arrow map is a regular epimorphism. Its kernel pair on arrows is an
internal category congruence, denoted
\[
\theta_q.
\]
\end{definition}

\begin{proposition}[Congruences and effective input quotients]
\label{prop:reiterman-congruences-effective-quotients}
In a Grothendieck topos, every internal category congruence
\[
\theta
\]
on \(F\) has an effective quotient internal category
\[
F/\theta.
\]
Consequently, effective input quotients of \(F\), up to isomorphism under
\(F\), are equivalently internal category congruences on \(F\).
\end{proposition}

\begin{proof}
Since \(\mathcal E\) is Barr-exact, every internal equivalence relation is
effective. Thus the arrow-object equivalence relation
\[
\theta\rightrightarrows F_1
\]
has a quotient
\[
F_1\twoheadrightarrow F_1/\theta.
\]
Because \(\theta\) lies over
\[
A_0\times A_0,
\]
the source and target maps descend to
\[
F_1/\theta\rightrightarrows A_0.
\]
Compatibility with identities and composition is exactly the condition that
the identity and composition maps of \(F\) descend through this quotient.

The descended structure maps satisfy the internal category axioms because
they do so after pullback along the regular epimorphism
\[
F_1\twoheadrightarrow F_1/\theta.
\]
Regular epimorphisms are stable under pullback and reflect equality of
descended maps. Hence \(F/\theta\) is an internal category and
\[
F\twoheadrightarrow F/\theta
\]
is an identity-on-objects regular quotient.

Conversely, the kernel pair of any identity-on-objects regular quotient of
internal categories is an internal category congruence. These two
constructions are inverse up to canonical isomorphism.
\end{proof}

\begin{definition}[Quotient order]
\label{def:reiterman-quotient-order}
For effective input quotients
\[
q:F\twoheadrightarrow C
\qquad\text{and}\qquad
p:F\twoheadrightarrow P,
\]
write
\[
q\preceq p
\]
if \(p\) is a further quotient of \(q\). Equivalently, there exists an
identity-on-objects regular quotient
\[
C\twoheadrightarrow P
\]
such that
\[
p=(C\twoheadrightarrow P)\circ q.
\]
In terms of congruences,
\[
q\preceq p
\quad\Longleftrightarrow\quad
\theta_q\leq\theta_p.
\]
\end{definition}

Thus larger congruences correspond to coarser quotients.

\section{Subdirect products of effective input quotients}
\label{sec:reiterman-subdirect-products-input-quotients}

\begin{definition}[Subdirect product of input quotients]
\label{def:reiterman-subdirect-product-input-quotients}
Let
\[
(q_i:F\twoheadrightarrow C_i)_{i\in I}
\]
be a small family of effective input quotients. Their \textbf{subdirect
product over \(F\)} is the image of the diagonal functor
\[
F\longrightarrow
\prod_{i\in I} C_i
\]
in the category of internal categories with object object \(A_0\).

Equivalently, it is the quotient
\[
F\twoheadrightarrow
F/\bigcap_{i\in I}\theta_{q_i}.
\]
\end{definition}

\begin{remark}[Empty subdirect product]
\label{rem:reiterman-empty-subdirect-product}
The empty subdirect product, when used, is the quotient of \(F\) by the
largest internal category congruence on parallel arrows. Equivalently, it is
the terminal identity-on-objects quotient of \(F\).
\end{remark}

\begin{proposition}[Kernel of a subdirect product]
\label{prop:reiterman-kernel-subdirect-product}
The kernel congruence of the subdirect product of
\[
(q_i:F\twoheadrightarrow C_i)_{i\in I}
\]
is
\[
\bigcap_{i\in I}\theta_{q_i}.
\]
\end{proposition}

\begin{proof}
Two arrows of \(F\) have the same image under the diagonal map
\[
F\to\prod_{i\in I}C_i
\]
if and only if they have the same image in each \(C_i\). This is precisely
membership in
\[
\bigcap_{i\in I}\theta_{q_i}.
\]
Taking the regular image of the diagonal therefore gives the quotient by this
intersection congruence.
\end{proof}

\begin{remark}
\label{rem:reiterman-small-subdirect-products}
For a regular cardinal \(\kappa\), \(<\kappa\)-subdirect products correspond
to intersections indexed by sets of cardinality \(<\kappa\).
\end{remark}

\section{Completed input categories and completed equations}
\label{sec:reiterman-completed-input-categories-equations}

Let
\[
\mathcal R
\]
be a small collection of effective input quotients of
\[
F=\mathbb A^{\mathrm{op}}.
\]
Assume that \(\mathcal R\) is closed under isomorphism of quotients under
\(F\) and under finite subdirect products. Then \(\mathcal R\), ordered by
further quotient maps, is cofiltered: for any two quotients, their subdirect
product maps to both of them.

The completion used in this section is a pro-internal category. It is used
only through its compatible projections to the quotients in \(\mathcal R\).

\begin{definition}[Completed input category]
\label{def:reiterman-completed-input-category}
The \textbf{\(\mathcal R\)-completed input category} is the pro-internal
category
\[
\widehat F^{\,\mathcal R}
\]
represented by the inverse system
\[
(q:F\twoheadrightarrow C)\in\mathcal R.
\]
Thus an arrow from a constant internal category \(T\) to
\[
\widehat F^{\,\mathcal R}
\]
is a compatible family of arrows
\[
T\to C
\]
as \(C\) ranges over quotients in \(\mathcal R\).
\end{definition}

\begin{definition}[Generalized completed input equation]
\label{def:reiterman-completed-input-equation}
A \textbf{generalized completed input equation} over \(\mathcal R\) is a pair
of generalized parallel arrows
\[
u,v:T\rightrightarrows \widehat F^{\,\mathcal R}_1
\]
in the pro-completion, over a common source-target map
\[
T\to A_0\times A_0.
\]
Equivalently, it is a pair of compatible families
\[
u_q,v_q:T\rightrightarrows C_{q,1},
\]
one for each quotient
\[
q:F\twoheadrightarrow C_q
\]
in \(\mathcal R\), such that \(u_q\) and \(v_q\) are parallel over
\[
A_0\times A_0
\]
and are compatible with further quotient maps.
\end{definition}

\begin{definition}[Satisfaction of a completed equation]
\label{def:reiterman-satisfaction-completed-equation}
Let
\[
q:F\twoheadrightarrow C
\]
belong to \(\mathcal R\). The quotient \(q\) \textbf{satisfies} the completed
input equation
\[
u=v
\]
if
\[
u_q=v_q:T\to C_1.
\]
\end{definition}

\begin{definition}[Models and equations]
\label{def:reiterman-models-equations}
For a class \(\Sigma\) of generalized completed input equations, write
\[
\operatorname{Mod}_{\mathcal R}(\Sigma)
\]
for the class of quotients in \(\mathcal R\) satisfying every equation in
\(\Sigma\).

For a class
\[
\mathcal Q\subseteq\mathcal R,
\]
write
\[
\operatorname{Eq}_{\mathcal R}(\mathcal Q)
\]
for the class of generalized completed input equations satisfied by every
quotient in \(\mathcal Q\).
\end{definition}

\begin{definition}[Ordinary generalized input equation]
\label{def:reiterman-ordinary-generalized-input-equation}
An \textbf{ordinary generalized input equation} is a pair of generalized
parallel input arrows
\[
u,v:T\rightrightarrows F_1
\]
over
\[
A_0\times A_0.
\]
It determines a completed equation over \(\mathcal R\) by applying every
quotient
\[
q:F\twoheadrightarrow C_q
\]
to obtain the compatible family
\[
q_1u,q_1v:T\rightrightarrows C_{q,1}.
\]
\end{definition}

\begin{lemma}[Ordinary equations detect congruence containment]
\label{lem:reiterman-ordinary-equations-detect-containment}
Let
\[
\mathcal Q\subseteq\mathcal R
\]
and let
\[
p:F\twoheadrightarrow P
\]
be a quotient in \(\mathcal R\). Then \(p\) satisfies every ordinary
generalized input equation valid in all quotients of \(\mathcal Q\) if and
only if
\[
\bigcap_{q\in\mathcal Q}\theta_q
\leq
\theta_p.
\]
\end{lemma}

\begin{proof}
The intersection
\[
\bigcap_{q\in\mathcal Q}\theta_q
\]
is a subobject of
\[
F_1\times_{A_0\times A_0}F_1.
\]
Containment in \(\theta_p\) says exactly that every generalized pair of
parallel arrows
\[
u,v:T\rightrightarrows F_1
\]
which is identified by every quotient \(q\in\mathcal Q\) is also identified
by \(p\). This is precisely satisfaction by \(p\) of all ordinary generalized
input equations valid in \(\mathcal Q\).
\end{proof}

\section{The complete equational theorem for effective input quotients}
\label{sec:reiterman-complete-equational-theorem}

The first quotient-equational theorem uses arbitrary small subdirect products.
It is the structural foundation for the compact Reiterman forms.

\begin{definition}[Complete-equational quotient class]
\label{def:reiterman-complete-equational-quotient-class}
Let
\[
\mathcal R
\]
be a small collection of effective input quotients closed under arbitrary
small subdirect products and further quotients. A class
\[
\mathcal Q\subseteq\mathcal R
\]
is called a \textbf{complete-equational quotient class} if it is closed under:
\begin{enumerate}
\item isomorphism of quotients under \(F\);
\item further quotients in \(\mathcal R\);
\item arbitrary small subdirect products.
\end{enumerate}
\end{definition}

\begin{theorem}[Complete equational theorem for effective input quotients]
\label{thm:reiterman-complete-equational-theorem}
Let
\[
\mathcal R
\]
be a small collection of effective input quotients of \(F\), closed under
arbitrary small subdirect products and further quotients. For a class
\[
\mathcal Q\subseteq\mathcal R,
\]
the following are equivalent:
\begin{enumerate}
\item \(\mathcal Q\) is a complete-equational quotient class.
\item There exists a class \(\Sigma\) of generalized completed input equations
over \(\mathcal R\) such that
\[
\mathcal Q=\operatorname{Mod}_{\mathcal R}(\Sigma).
\]
\end{enumerate}
Moreover, one may take
\[
\Sigma=\operatorname{Eq}_{\mathcal R}(\mathcal Q).
\]
\end{theorem}

\begin{proof}
Suppose first that
\[
\mathcal Q=\operatorname{Mod}_{\mathcal R}(\Sigma).
\]
Satisfaction of completed equations is invariant under isomorphism of
quotients under \(F\). It is preserved by further quotients because equality
descends along the projection to a quotient.

It is preserved by arbitrary small subdirect products as follows. Let
\[
r:F\twoheadrightarrow R
\]
be the subdirect product of quotients
\[
q_i:F\twoheadrightarrow C_i
\]
which all satisfy a completed equation
\[
u=v.
\]
The quotient \(R\) is a subobject of
\[
\prod_i C_i
\]
over the common object object \(A_0\). The projections
\[
R\to C_i
\]
are jointly monic on arrows. Since
\[
u_{q_i}=v_{q_i}
\]
for every \(i\), joint monicity gives
\[
u_r=v_r.
\]
Thus \(r\) satisfies the equation. Therefore \(\mathcal Q\) is closed under
isomorphism, further quotients, and arbitrary small subdirect products.

Conversely, suppose \(\mathcal Q\) is closed under those operations. We show
\[
\mathcal Q
=
\operatorname{Mod}_{\mathcal R}
(\operatorname{Eq}_{\mathcal R}(\mathcal Q)).
\]
The inclusion
\[
\mathcal Q
\subseteq
\operatorname{Mod}_{\mathcal R}
(\operatorname{Eq}_{\mathcal R}(\mathcal Q))
\]
is immediate.

For the converse, let
\[
p:F\twoheadrightarrow P
\]
be a quotient in \(\mathcal R\) satisfying every generalized completed input
equation valid in all quotients of \(\mathcal Q\). In particular, \(p\)
satisfies every ordinary generalized input equation valid in \(\mathcal Q\).
By Lemma~\ref{lem:reiterman-ordinary-equations-detect-containment},
\[
\bigcap_{q\in\mathcal Q}\theta_q
\leq
\theta_p.
\]

Let
\[
r:F\twoheadrightarrow R
\]
be the subdirect product of all quotients in \(\mathcal Q\), using
Remark~\ref{rem:reiterman-empty-subdirect-product} when \(\mathcal Q\) is
empty. By Proposition~\ref{prop:reiterman-kernel-subdirect-product},
\[
\theta_r
=
\bigcap_{q\in\mathcal Q}\theta_q.
\]
Since \(\mathcal Q\) is closed under arbitrary small subdirect products,
\[
r\in\mathcal Q.
\]
The inequality
\[
\theta_r\leq\theta_p
\]
says exactly that \(p\) is a further quotient of \(r\). Since \(\mathcal Q\)
is closed under further quotients in \(\mathcal R\), it follows that
\[
p\in\mathcal Q.
\]
Thus
\[
\operatorname{Mod}_{\mathcal R}
(\operatorname{Eq}_{\mathcal R}(\mathcal Q))
\subseteq
\mathcal Q.
\]
This proves the theorem.
\end{proof}

\section{Intersection-compact input quotients}
\label{sec:reiterman-intersection-compact-input-quotients}

The abstract Reiterman form is obtained by replacing arbitrary subdirect
products with small subdirect products. The required finite-extraction
principle is intersection compactness.

Let
\[
\kappa
\]
be a regular cardinal.

\begin{definition}[\(\kappa\)-intersection-compact quotient]
\label{def:reiterman-kappa-intersection-compact-quotient}
An effective input quotient
\[
p:F\twoheadrightarrow P
\]
is called \textbf{\(\kappa\)-intersection-compact} if its kernel congruence
\[
\theta_p
\]
satisfies the following property.

Whenever a small family of internal category congruences
\[
(\theta_i)_{i\in I}
\]
on \(F\) satisfies
\[
\bigcap_{i\in I}\theta_i\leq\theta_p,
\]
there exists a subset
\[
J\subseteq I
\]
with
\[
|J|<\kappa
\]
such that
\[
\bigcap_{j\in J}\theta_j\leq\theta_p.
\]
\end{definition}

When
\[
\kappa=\omega,
\]
we simply say \textbf{intersection-compact}.

\begin{proposition}[Small subdirect products preserve compactness]
\label{prop:reiterman-small-subdirect-products-preserve-compactness}
Let
\[
(q_j:F\twoheadrightarrow C_j)_{j\in J}
\]
be a family of \(\kappa\)-intersection-compact effective input quotients with
\[
|J|<\kappa.
\]
Then their subdirect product is again \(\kappa\)-intersection-compact.
\end{proposition}

\begin{proof}
Let
\[
q:F\twoheadrightarrow C
\]
be the subdirect product. Its kernel congruence is
\[
\theta_q=\bigcap_{j\in J}\theta_{q_j}.
\]
Suppose
\[
\bigcap_{i\in I}\phi_i\leq\theta_q.
\]
Then for every \(j\in J\),
\[
\bigcap_{i\in I}\phi_i\leq\theta_{q_j}.
\]
Since \(q_j\) is \(\kappa\)-intersection-compact, there exists
\[
I_j\subseteq I
\]
with
\[
|I_j|<\kappa
\]
and
\[
\bigcap_{i\in I_j}\phi_i\leq\theta_{q_j}.
\]
Let
\[
I_0:=\bigcup_{j\in J}I_j.
\]
Since \(\kappa\) is regular and \(|J|<\kappa\), one has
\[
|I_0|<\kappa.
\]
Moreover,
\[
\bigcap_{i\in I_0}\phi_i
\leq
\theta_{q_j}
\]
for every \(j\in J\). Hence
\[
\bigcap_{i\in I_0}\phi_i
\leq
\bigcap_{j\in J}\theta_{q_j}
=
\theta_q.
\]
Thus \(q\) is \(\kappa\)-intersection-compact.
\end{proof}

\begin{definition}[\(\kappa\)-compact input quotient collection]
\label{def:reiterman-kappa-compact-input-quotient-collection}
A \textbf{\(\kappa\)-compact input quotient collection} for \(F\) is a small
collection
\[
\mathcal R_\kappa
\]
of \(\kappa\)-intersection-compact effective input quotients of \(F\), closed
under:
\begin{enumerate}
\item isomorphism of quotients under \(F\);
\item further quotients that are \(\kappa\)-intersection-compact;
\item \(<\kappa\)-subdirect products.
\end{enumerate}
\end{definition}

\begin{remark}[Intersection compactness versus finiteness]
\label{rem:reiterman-intersection-compactness-versus-finiteness}
Intersection compactness is an abstract meet-compactness condition on kernel
congruences. It is not the same as finiteness of the quotient. The
finite/profinite Reiterman theorem below uses a different compactness
principle: compactness of profinite inverse limits and finite determinacy of
continuous maps into finite discrete quotients.
\end{remark}

\section{The intersection-compact Reiterman theorem}
\label{sec:reiterman-intersection-compact}

\begin{definition}[\(\kappa\)-compact algebraic quotient class]
\label{def:reiterman-kappa-compact-algebraic-quotient-class}
Let
\[
\mathcal R_\kappa
\]
be a \(\kappa\)-compact input quotient collection. A class
\[
\mathcal Q\subseteq\mathcal R_\kappa
\]
is called a \textbf{\(\kappa\)-compact algebraic quotient class} if it is
closed under:
\begin{enumerate}
\item isomorphism of quotients under \(F\);
\item further quotients in \(\mathcal R_\kappa\);
\item \(<\kappa\)-subdirect products.
\end{enumerate}
\end{definition}

\begin{theorem}[Intersection-compact Reiterman theorem]
\label{thm:reiterman-intersection-compact}
Let
\[
\mathcal R_\kappa
\]
be a \(\kappa\)-compact input quotient collection for
\[
F=\mathbb A^{\mathrm{op}}.
\]
For a class
\[
\mathcal Q\subseteq\mathcal R_\kappa,
\]
the following are equivalent:
\begin{enumerate}
\item \(\mathcal Q\) is a \(\kappa\)-compact algebraic quotient class.
\item There exists a class \(\Sigma\) of generalized completed input equations
over \(\mathcal R_\kappa\) such that
\[
\mathcal Q=\operatorname{Mod}_{\mathcal R_\kappa}(\Sigma).
\]
\end{enumerate}
Moreover, one may take
\[
\Sigma=\operatorname{Eq}_{\mathcal R_\kappa}(\mathcal Q).
\]
\end{theorem}

\begin{proof}
Suppose first that
\[
\mathcal Q=\operatorname{Mod}_{\mathcal R_\kappa}(\Sigma).
\]
As in Theorem~\ref{thm:reiterman-complete-equational-theorem}, satisfaction
of completed input equations is invariant under isomorphism, preserved by
further quotients, and preserved by subdirect products. Hence it is preserved
by \(<\kappa\)-subdirect products. Therefore \(\mathcal Q\) is a
\(\kappa\)-compact algebraic quotient class.

Conversely, suppose that \(\mathcal Q\) is closed under isomorphism, further
quotients in \(\mathcal R_\kappa\), and \(<\kappa\)-subdirect products. We
prove
\[
\mathcal Q
=
\operatorname{Mod}_{\mathcal R_\kappa}
(\operatorname{Eq}_{\mathcal R_\kappa}(\mathcal Q)).
\]
The inclusion
\[
\mathcal Q
\subseteq
\operatorname{Mod}_{\mathcal R_\kappa}
(\operatorname{Eq}_{\mathcal R_\kappa}(\mathcal Q))
\]
is immediate.

For the converse, let
\[
p:F\twoheadrightarrow P
\]
be a quotient in \(\mathcal R_\kappa\) satisfying every generalized completed
input equation valid in all quotients of \(\mathcal Q\). In particular, \(p\)
satisfies every ordinary generalized input equation valid in \(\mathcal Q\).
By Lemma~\ref{lem:reiterman-ordinary-equations-detect-containment},
\[
\bigcap_{q\in\mathcal Q}\theta_q
\leq
\theta_p.
\]
Since
\[
p\in\mathcal R_\kappa,
\]
the quotient \(p\) is \(\kappa\)-intersection-compact. Hence there is a
subclass
\[
\mathcal Q_0\subseteq\mathcal Q
\]
with
\[
|\mathcal Q_0|<\kappa
\]
such that
\[
\bigcap_{q\in\mathcal Q_0}\theta_q
\leq
\theta_p.
\]
Let
\[
r:F\twoheadrightarrow R
\]
be the subdirect product of the quotients in \(\mathcal Q_0\). Then
\[
\theta_r
=
\bigcap_{q\in\mathcal Q_0}\theta_q.
\]
Because \(\mathcal Q\) is closed under \(<\kappa\)-subdirect products,
\[
r\in\mathcal Q.
\]
The inequality
\[
\theta_r\leq\theta_p
\]
says that \(p\) is a further quotient of \(r\). Since \(\mathcal Q\) is
closed under further quotients in \(\mathcal R_\kappa\), it follows that
\[
p\in\mathcal Q.
\]
Thus
\[
\operatorname{Mod}_{\mathcal R_\kappa}
(\operatorname{Eq}_{\mathcal R_\kappa}(\mathcal Q))
\subseteq
\mathcal Q.
\]
This proves the theorem.
\end{proof}

\section{Reiterman-admissible interface translations}
\label{sec:reiterman-admissible-interface-translations}

We now pass from the local quotient theory at a fixed interface
\[
\Xi=(\mathbb A,\mathbb B)
\]
to a global theory over varying interfaces. The state-level part of an
interface translation was isolated in the preceding chapter. The quotient
theory requires additional algebraic data.

\begin{definition}[Mealy interface]
\label{def:reiterman-mealy-interface}
A \textbf{Mealy interface} is a pair
\[
\Xi=(\mathbb A_\Xi,\mathbb B_\Xi)
\]
of internal categories in the ambient topos. We write
\[
\mathsf{Beh}_\Xi
:=
\mathsf{Beh}_{\mathbb A_\Xi,\mathbb B_\Xi},
\]
\[
Z_\Xi:=
(A_\Xi)_0\times(B_\Xi)_0,
\]
\[
F_\Xi:=\mathbb A_\Xi^{\mathrm{op}},
\]
and
\[
\mathsf D_\Xi
\]
for the residual closure operator on
\[
\operatorname{Sub}_{Z_\Xi}(\mathsf{Beh}_\Xi).
\]
\end{definition}

\begin{definition}[Effective quotient poset of an interface]
\label{def:reiterman-effective-quotient-poset-interface}
For an interface \(\Xi\), define
\[
\mathsf{Quot}_\Xi
\]
to be the poset of effective input quotients of
\[
F_\Xi=\mathbb A_\Xi^{\mathrm{op}},
\]
ordered by further quotient.
\end{definition}

\begin{definition}[Reiterman-admissible interface translation]
\label{def:reiterman-admissible-interface-translation}
Let
\[
\Xi'=(\mathbb A',\mathbb B')
\qquad\text{and}\qquad
\Xi=(\mathbb A,\mathbb B)
\]
be interfaces. A \textbf{Reiterman-admissible interface translation}
\[
\Gamma:\Xi'\to\Xi
\]
consists of the following data.

\begin{enumerate}
\item A state-admissible translation on behaviour specifications
\[
\Gamma_{\mathrm{beh}}^\ast:
\operatorname{Sub}_{Z_\Xi}(\mathsf{Beh}_\Xi)
\to
\operatorname{Sub}_{Z_{\Xi'}}(\mathsf{Beh}_{\Xi'})
\]
compatible with residual closure and the chosen behaviour-operation
doctrines.

\item An internal functor on input categories
\[
\gamma_\Gamma:F_{\Xi'}\to F_\Xi.
\]

\item A reindexing operation on effective input quotients
\[
\Gamma_{\mathrm{alg}}^\ast:
\mathsf{Quot}_{\Xi}
\to
\mathsf{Quot}_{\Xi'}
\]
defined as follows. If
\[
q:F_\Xi\twoheadrightarrow C
\]
has kernel congruence \(\theta_q\), then
\[
\Gamma_{\mathrm{alg}}^\ast(q)
\]
is the effective quotient of \(F_{\Xi'}\) by the inverse-image congruence
\[
\gamma_\Gamma^{-1}(\theta_q)
\hookrightarrow
(F_{\Xi'})_1
\times_{(A_{\Xi'})_0\times(A_{\Xi'})_0}
(F_{\Xi'})_1.
\]

\item Compatibility with algebraic recognizers: for every closed recognizer
\[
D\hookrightarrow\mathsf{Beh}_{\Xi},
\]
there is a canonical isomorphism of effective input quotients of \(F_{\Xi'}\)
\[
\Gamma_{\mathrm{alg}}^\ast
\bigl(
\operatorname{Alg}_{\Xi}(D)
\bigr)
\cong
\operatorname{Alg}_{\Xi'}
\bigl(
\Gamma_{\mathrm{st}}^\ast D
\bigr),
\]
where
\[
\Gamma_{\mathrm{st}}^\ast D
:=
\mathsf D_{\Xi'}(\Gamma_{\mathrm{beh}}^\ast D).
\]
\end{enumerate}
\end{definition}

\begin{proposition}[Syntactic algebraic recognizers commute with translations]
\label{prop:reiterman-synalg-commutes-with-translations}
Let
\[
\Gamma:\Xi'\to\Xi
\]
be a Reiterman-admissible interface translation. Then for every behaviour
specification
\[
S\in\operatorname{Sub}_{Z_\Xi}(\mathsf{Beh}_\Xi),
\]
there is a canonical isomorphism of effective input quotients of \(F_{\Xi'}\)
\[
\Gamma_{\mathrm{alg}}^\ast
\bigl(
\operatorname{SynAlg}_{\Xi}(S)
\bigr)
\cong
\operatorname{SynAlg}_{\Xi'}
\bigl(
\Gamma_{\mathrm{beh}}^\ast S
\bigr).
\]
\end{proposition}

\begin{proof}
By definition,
\[
\operatorname{SynAlg}_{\Xi}(S)
=
\operatorname{Alg}_{\Xi}(\mathsf D_\Xi(S)).
\]
Therefore
\[
\Gamma_{\mathrm{alg}}^\ast
\bigl(
\operatorname{SynAlg}_{\Xi}(S)
\bigr)
=
\Gamma_{\mathrm{alg}}^\ast
\bigl(
\operatorname{Alg}_{\Xi}(\mathsf D_\Xi(S))
\bigr).
\]
By algebraic compatibility in
Definition~\ref{def:reiterman-admissible-interface-translation}, this is
canonically isomorphic to
\[
\operatorname{Alg}_{\Xi'}
\bigl(
\Gamma_{\mathrm{st}}^\ast(\mathsf D_\Xi(S))
\bigr).
\]
By definition of \(\Gamma_{\mathrm{st}}^\ast\),
\[
\Gamma_{\mathrm{st}}^\ast(\mathsf D_\Xi(S))
=
\mathsf D_{\Xi'}
\bigl(
\Gamma_{\mathrm{beh}}^\ast(\mathsf D_\Xi(S))
\bigr).
\]
By residual compatibility,
\[
\Gamma_{\mathrm{beh}}^\ast(\mathsf D_\Xi(S))
=
\mathsf D_{\Xi'}(\Gamma_{\mathrm{beh}}^\ast S).
\]
Using idempotence of \(\mathsf D_{\Xi'}\), we get
\[
\Gamma_{\mathrm{st}}^\ast(\mathsf D_\Xi(S))
=
\mathsf D_{\Xi'}(\Gamma_{\mathrm{beh}}^\ast S).
\]
Hence
\[
\Gamma_{\mathrm{alg}}^\ast
\bigl(
\operatorname{SynAlg}_{\Xi}(S)
\bigr)
\cong
\operatorname{Alg}_{\Xi'}
\bigl(
\mathsf D_{\Xi'}(\Gamma_{\mathrm{beh}}^\ast S)
\bigr)
=
\operatorname{SynAlg}_{\Xi'}
\bigl(
\Gamma_{\mathrm{beh}}^\ast S
\bigr).
\]
\end{proof}

\section{Global algebraic quotient theories}
\label{sec:reiterman-global-algebraic-quotient-theories}

Let
\[
\mathfrak T
\]
be a category of Mealy interfaces and Reiterman-admissible translations.

\begin{definition}[Global algebraic quotient collection]
\label{def:reiterman-global-algebraic-quotient-collection}
A \textbf{global algebraic quotient collection}
\[
\mathcal R=(\mathcal R_\Xi)_\Xi
\]
on \(\mathfrak T\) assigns to every interface \(\Xi\) a small collection
\[
\mathcal R_\Xi\subseteq\mathsf{Quot}_\Xi
\]
of effective input quotients such that:
\begin{enumerate}
\item each \(\mathcal R_\Xi\) is closed under isomorphism of quotients under
\(F_\Xi\);
\item each \(\mathcal R_\Xi\) is closed under the local quotient operations
specified in context;
\item for every Reiterman-admissible translation
\[
\Gamma:\Xi'\to\Xi
\]
and every
\[
q\in\mathcal R_\Xi,
\]
one has
\[
\Gamma_{\mathrm{alg}}^\ast q\in\mathcal R_{\Xi'}.
\]
\end{enumerate}
\end{definition}

\begin{definition}[Global algebraic quotient class]
\label{def:reiterman-global-algebraic-quotient-class}
Let
\[
\mathcal R=(\mathcal R_\Xi)_\Xi
\]
be a global algebraic quotient collection. A \textbf{global algebraic
quotient class}
\[
\mathcal Q=(\mathcal Q_\Xi)_\Xi
\]
with
\[
\mathcal Q_\Xi\subseteq\mathcal R_\Xi
\]
is \textbf{translation-stable} if for every Reiterman-admissible translation
\[
\Gamma:\Xi'\to\Xi
\]
and every
\[
q\in\mathcal Q_\Xi,
\]
one has
\[
\Gamma_{\mathrm{alg}}^\ast q\in\mathcal Q_{\Xi'}.
\]
\end{definition}

\begin{proposition}[Syntactic quotients of global varieties are translation-stable]
\label{prop:reiterman-syntactic-quotients-global-varieties}
Let
\[
\mathcal V=(\mathcal V_\Xi)_\Xi
\]
be a global behaviour variety in the sense of the state Eilenberg theorem,
and assume the interface translations are Reiterman-admissible. Define
\[
\operatorname{SynAlg}(\mathcal V)_\Xi
:=
\left\{
\operatorname{SynAlg}_{\Xi}(S)
\ \middle|\
S\in\mathcal V_\Xi
\right\}.
\]
Then
\[
\operatorname{SynAlg}(\mathcal V)
\]
is stable under Reiterman-admissible translations.
\end{proposition}

\begin{proof}
Let
\[
\Gamma:\Xi'\to\Xi
\]
be Reiterman-admissible and let
\[
\operatorname{SynAlg}_{\Xi}(S)
\]
belong to
\[
\operatorname{SynAlg}(\mathcal V)_\Xi.
\]
By Proposition~\ref{prop:reiterman-synalg-commutes-with-translations},
\[
\Gamma_{\mathrm{alg}}^\ast
\bigl(
\operatorname{SynAlg}_{\Xi}(S)
\bigr)
\cong
\operatorname{SynAlg}_{\Xi'}
\bigl(
\Gamma_{\mathrm{beh}}^\ast S
\bigr).
\]
Since
\[
S\in\mathcal V_\Xi
\]
and \(\mathcal V\) is global,
\[
\Gamma_{\mathrm{beh}}^\ast S\in\mathcal V_{\Xi'}.
\]
Therefore the right-hand side belongs to
\[
\operatorname{SynAlg}(\mathcal V)_{\Xi'}.
\]
\end{proof}

\section{Global completed input equations}
\label{sec:reiterman-global-completed-input-equations}

Let
\[
\mathcal R=(\mathcal R_\Xi)_\Xi
\]
be a global algebraic quotient collection. For each interface \(\Xi\), form
the completed input category
\[
\widehat F_\Xi^{\,\mathcal R_\Xi}
\]
from the local quotient collection \(\mathcal R_\Xi\).

\begin{definition}[Global completed equation datum]
\label{def:reiterman-global-completed-equation-datum}
A \textbf{global completed equation datum} on \(\mathcal R\) assigns:
\begin{enumerate}
\item to every interface \(\Xi\), the class of generalized completed input
equations over
\[
\widehat F_\Xi^{\,\mathcal R_\Xi};
\]
\item to every Reiterman-admissible translation
\[
\Gamma:\Xi'\to\Xi,
\]
a reindexing operation on completed equations
\[
\Gamma^\sharp:
\operatorname{Eqn}_{\mathcal R_\Xi}(\Xi)
\to
\operatorname{Eqn}_{\mathcal R_{\Xi'}}(\Xi')
\]
such that satisfaction is compatible with quotient reindexing:
\[
q\models e
\quad\Longrightarrow\quad
\Gamma_{\mathrm{alg}}^\ast q\models \Gamma^\sharp(e)
\]
for every
\[
q\in\mathcal R_\Xi.
\]
\end{enumerate}
\end{definition}

In concrete profinite situations, \(\Gamma^\sharp\) is induced by the
continuous extension of the input functor
\[
\gamma_\Gamma:F_{\Xi'}\to F_\Xi
\]
to completed input categories.

\begin{definition}[Global completed equational theory]
\label{def:reiterman-global-completed-equational-theory}
A \textbf{global completed equational theory}
\[
\Sigma=(\Sigma_\Xi)_\Xi
\]
on \(\mathcal R\) assigns to every interface \(\Xi\) a class
\[
\Sigma_\Xi
\]
of generalized completed input equations over
\[
\widehat F_\Xi^{\,\mathcal R_\Xi}
\]
and is required to be \textbf{model-stable under admissible translation}:
for every Reiterman-admissible translation
\[
\Gamma:\Xi'\to\Xi
\]
and every quotient
\[
q\in\mathcal R_\Xi,
\]
if
\[
q\models \Sigma_\Xi,
\]
then
\[
\Gamma_{\mathrm{alg}}^\ast q
\models
\Sigma_{\Xi'}.
\]
\end{definition}

This model-stability condition is the global compatibility needed for
Reiterman theory. Closure merely under syntactic reindexing of displayed
equations is not by itself sufficient unless the theory is generated by those
reindexed equations or the additional model-stability condition is imposed.

\begin{definition}[Models of a global completed equational theory]
\label{def:reiterman-models-global-equational-theory}
For a global completed equational theory
\[
\Sigma=(\Sigma_\Xi)_\Xi,
\]
define
\[
\operatorname{Mod}_{\mathcal R}(\Sigma)_\Xi
:=
\left\{
q\in\mathcal R_\Xi
\ \middle|\
q\models e\text{ for every }e\in\Sigma_\Xi
\right\}.
\]
\end{definition}

\begin{proposition}[Models are translation-stable]
\label{prop:reiterman-global-models-translation-stable}
If
\[
\Sigma
\]
is a global completed equational theory, then
\[
\operatorname{Mod}_{\mathcal R}(\Sigma)
\]
is stable under Reiterman-admissible translations.
\end{proposition}

\begin{proof}
This is exactly the model-stability condition in
Definition~\ref{def:reiterman-global-completed-equational-theory}. If
\[
q\in\operatorname{Mod}_{\mathcal R}(\Sigma)_\Xi
\]
and
\[
\Gamma:\Xi'\to\Xi
\]
is Reiterman-admissible, then
\[
q\models\Sigma_\Xi.
\]
Therefore
\[
\Gamma_{\mathrm{alg}}^\ast q
\models
\Sigma_{\Xi'},
\]
which says precisely that
\[
\Gamma_{\mathrm{alg}}^\ast q
\in
\operatorname{Mod}_{\mathcal R}(\Sigma)_{\Xi'}.
\]
\end{proof}

\section{Global Reiterman theorem: intersection-compact form}
\label{sec:reiterman-global-intersection-compact}

\begin{definition}[Global \(\kappa\)-compact quotient collection]
\label{def:reiterman-global-kappa-compact-quotient-collection}
A global algebraic quotient collection
\[
\mathcal R_\kappa=(\mathcal R_{\kappa,\Xi})_\Xi
\]
is \textbf{\(\kappa\)-compact} if each local collection
\[
\mathcal R_{\kappa,\Xi}
\]
is a \(\kappa\)-compact input quotient collection and Reiterman-admissible
translation reindexing preserves the collection:
\[
q\in\mathcal R_{\kappa,\Xi}
\quad\Longrightarrow\quad
\Gamma_{\mathrm{alg}}^\ast q\in\mathcal R_{\kappa,\Xi'}.
\]
\end{definition}

\begin{definition}[Global \(\kappa\)-compact algebraic quotient class]
\label{def:reiterman-global-kappa-compact-algebraic-class}
Let
\[
\mathcal R_\kappa
\]
be a global \(\kappa\)-compact quotient collection. A class
\[
\mathcal Q=(\mathcal Q_\Xi)_\Xi,
\qquad
\mathcal Q_\Xi\subseteq\mathcal R_{\kappa,\Xi},
\]
is a \textbf{global \(\kappa\)-compact algebraic quotient class} if:
\begin{enumerate}
\item each \(\mathcal Q_\Xi\) is closed under isomorphism of quotients under
\(F_\Xi\);
\item each \(\mathcal Q_\Xi\) is closed under further quotients in
\(\mathcal R_{\kappa,\Xi}\);
\item each \(\mathcal Q_\Xi\) is closed under \(<\kappa\)-subdirect products;
\item \(\mathcal Q\) is stable under Reiterman-admissible translations.
\end{enumerate}
\end{definition}

\begin{theorem}[Global Reiterman theorem, intersection-compact form]
\label{thm:reiterman-global-intersection-compact}
Let
\[
\mathcal R_\kappa
\]
be a global \(\kappa\)-compact quotient collection equipped with a compatible
global completed equation datum. For a family
\[
\mathcal Q=(\mathcal Q_\Xi)_\Xi,
\qquad
\mathcal Q_\Xi\subseteq\mathcal R_{\kappa,\Xi},
\]
the following are equivalent:
\begin{enumerate}
\item \(\mathcal Q\) is a global \(\kappa\)-compact algebraic quotient class.
\item There exists a global completed equational theory
\[
\Sigma=(\Sigma_\Xi)_\Xi
\]
such that
\[
\mathcal Q_\Xi
=
\operatorname{Mod}_{\mathcal R_{\kappa,\Xi}}(\Sigma_\Xi)
\]
for every interface \(\Xi\).
\end{enumerate}
Moreover, one may take
\[
\Sigma_\Xi
=
\operatorname{Eq}_{\mathcal R_{\kappa,\Xi}}(\mathcal Q_\Xi),
\]
with global model-stability supplied by the local theorem and translation
stability of \(\mathcal Q\).
\end{theorem}

\begin{proof}
The local equivalence at each interface \(\Xi\) is
Theorem~\ref{thm:reiterman-intersection-compact}. It remains to check
compatibility with Reiterman-admissible translations.

Suppose first that
\[
\mathcal Q=\operatorname{Mod}_{\mathcal R}(\Sigma)
\]
for a global completed equational theory \(\Sigma\). Local closure under
isomorphism, further quotients, and \(<\kappa\)-subdirect products follows
objectwise from Theorem~\ref{thm:reiterman-intersection-compact}.
Translation stability follows from
Proposition~\ref{prop:reiterman-global-models-translation-stable}.

Conversely, suppose \(\mathcal Q\) is a global \(\kappa\)-compact algebraic
quotient class. Objectwise, define
\[
\Sigma_\Xi
:=
\operatorname{Eq}_{\mathcal R_{\kappa,\Xi}}(\mathcal Q_\Xi).
\]
By the local theorem,
\[
\mathcal Q_\Xi
=
\operatorname{Mod}_{\mathcal R_{\kappa,\Xi}}(\Sigma_\Xi).
\]
We show that the family \(\Sigma=(\Sigma_\Xi)_\Xi\) is model-stable under
admissible translations. Let
\[
q\in\mathcal R_{\kappa,\Xi}
\]
satisfy
\[
q\models\Sigma_\Xi.
\]
By the local theorem,
\[
q\in\mathcal Q_\Xi.
\]
If
\[
\Gamma:\Xi'\to\Xi
\]
is Reiterman-admissible, translation stability of \(\mathcal Q\) gives
\[
\Gamma_{\mathrm{alg}}^\ast q\in\mathcal Q_{\Xi'}.
\]
Again by the local theorem,
\[
\Gamma_{\mathrm{alg}}^\ast q
\models
\Sigma_{\Xi'}.
\]
Thus \(\Sigma\) is a global completed equational theory, and its model class
is exactly \(\mathcal Q\).
\end{proof}

\section{Finite/profinite input quotients in Set}
\label{sec:reiterman-finite-profinite-input-quotients}

The preceding Reiterman theorem uses intersection compactness. In classical
finite recognition, the relevant compactness principle is different. It is
the compactness of profinite completions and the finite determinacy of
continuous maps into finite discrete quotients.

For this section, assume
\[
\mathcal E=\mathbf{Set}.
\]
Let
\[
F
\]
be a small category, viewed as a Set-internal category, and assume that its
object set
\[
F_0
\]
is finite. We work with identity-on-objects quotients of \(F\). Thus a finite
effective input quotient
\[
q:F\twoheadrightarrow C
\]
has the same finite object set as \(F\), and finite arrow set \(C_1\).

\begin{definition}[Finite quotient collection]
\label{def:reiterman-finite-quotient-collection}
A \textbf{finite quotient collection}
\[
\mathcal R^{\mathrm{fin}}
\]
for \(F\) is a small cofiltered collection of finite effective input
quotients of
\[
F
\]
closed under isomorphism, further finite quotients, and finite subdirect
products.
\end{definition}

\begin{definition}[Profinite completion]
\label{def:reiterman-profinite-completion}
The \textbf{profinite completion of \(F\) with respect to
\(\mathcal R^{\mathrm{fin}}\)} is the inverse limit
\[
\widehat F^{\,\mathrm{fin}}
:=
\varprojlim_{q:F\twoheadrightarrow C\in\mathcal R^{\mathrm{fin}}} C
\]
in the category of topological internal categories, where each finite
quotient \(C\) carries the discrete topology.
\end{definition}

Thus the arrow space of
\[
\widehat F^{\,\mathrm{fin}}
\]
is compact, Hausdorff, and totally disconnected. Since the object set is
finite and the quotients are identity-on-objects, the object space remains a
finite discrete space.

\begin{definition}[Profinite completed equation]
\label{def:reiterman-profinite-completed-equation}
A \textbf{profinite completed input equation} is a pair of parallel arrows
\[
u,v:T\rightrightarrows \widehat F^{\,\mathrm{fin}}_1
\]
where \(T\) is a finite discrete parameter set, or more generally a compact
totally disconnected parameter space, and the maps are continuous and have a
common source-target map.
\end{definition}

\begin{definition}[Satisfaction of profinite equations]
\label{def:reiterman-profinite-satisfaction}
A finite quotient
\[
q:F\twoheadrightarrow C
\]
satisfies the profinite equation
\[
u=v
\]
if the composites
\[
T\rightrightarrows
\widehat F^{\,\mathrm{fin}}_1
\to
C_1
\]
are equal.
\end{definition}

\begin{lemma}[Finite determinacy]
\label{lem:reiterman-finite-determinacy}
Let
\[
X=\varprojlim_{i\in I}X_i
\]
be an inverse limit of finite discrete sets over a cofiltered diagram. Let
\[
Y
\]
be a finite discrete set. Every continuous map
\[
f:X\to Y
\]
factors through a finite subproduct, equivalently through some finite
subdirect quotient determined by finitely many coordinates.
\end{lemma}

\begin{proof}
For each
\[
y\in Y,
\]
the fibre
\[
f^{-1}(y)
\]
is clopen in \(X\). The topology on \(X\) has a basis consisting of inverse
images of subsets of finite coordinate projections. Since \(Y\) is finite,
the clopen partition
\[
X=\coprod_{y\in Y}f^{-1}(y)
\]
is finite. Each part is compact, hence covered by finitely many basic
clopens. Taking the finitely many coordinates appearing in all these basic
clopens, one obtains a finite coordinate projection
\[
X\to X_J
\]
through which \(f\) factors.
\end{proof}

\begin{proposition}[Finite maps factor through finite subdirect products]
\label{prop:reiterman-finite-maps-factor-subdirect}
Let
\[
\mathcal Q\subseteq\mathcal R^{\mathrm{fin}}
\]
be a class of finite quotients, and let
\[
\widehat F_{\mathcal Q}
:=
\varprojlim_{q\in\mathcal Q}C_q
\]
be the pro-\(\mathcal Q\) completion. Every continuous quotient functor
\[
\widehat F_{\mathcal Q}\to P
\]
to a finite discrete category \(P\) factors through a finite subdirect
product of quotients in \(\mathcal Q\).
\end{proposition}

\begin{proof}
Apply Lemma~\ref{lem:reiterman-finite-determinacy} to the arrow map. The
finitely many coordinates used define a finite subdirect product
\[
R
\]
of quotients in \(\mathcal Q\). Compatibility with source, target, identities,
and composition follows because the original map is a continuous functor and
the factorization is through coordinate quotients of the profinite internal
category. Thus the functor factors through
\[
R.
\]
\end{proof}

\section{The finite/profinite Reiterman theorem}
\label{sec:reiterman-finite-profinite-theorem}

\begin{theorem}[Finite/profinite Reiterman theorem]
\label{thm:reiterman-finite-profinite}
Let
\[
\mathcal R^{\mathrm{fin}}
\]
be a finite quotient collection for a fixed input category \(F\) with finite
object set. For a class
\[
\mathcal Q\subseteq\mathcal R^{\mathrm{fin}},
\]
the following are equivalent:
\begin{enumerate}
\item \(\mathcal Q\) is closed under isomorphism, further finite quotients,
and finite subdirect products.
\item There exists a class \(\Sigma\) of profinite completed input equations
such that
\[
\mathcal Q=\operatorname{Mod}_{\mathcal R^{\mathrm{fin}}}(\Sigma).
\]
\end{enumerate}
Moreover, one may take
\[
\Sigma
\]
to be the class of all profinite completed equations valid in every quotient
of \(\mathcal Q\).
\end{theorem}

\begin{proof}
Suppose first that
\[
\mathcal Q=\operatorname{Mod}_{\mathcal R^{\mathrm{fin}}}(\Sigma).
\]
Satisfaction of profinite equations is invariant under isomorphism of
quotients and is preserved by further quotients. It is also preserved by
finite subdirect products because finite subdirect products are subobjects of
finite products with jointly monic projections, and equality can be checked
coordinatewise. Hence \(\mathcal Q\) has the stated closure properties.

Conversely, suppose that \(\mathcal Q\) is closed under isomorphism, further
finite quotients, and finite subdirect products. Let
\[
p:F\twoheadrightarrow P
\]
be a finite quotient satisfying every profinite completed equation valid in
\(\mathcal Q\). We prove that
\[
p\in\mathcal Q.
\]

Let
\[
\widehat F
\]
be the profinite completion with respect to
\[
\mathcal R^{\mathrm{fin}},
\]
and let
\[
\widehat F_{\mathcal Q}
\]
be the pro-\(\mathcal Q\) completion. The canonical projections induce a
continuous functor
\[
\pi_{\mathcal Q}:\widehat F\to\widehat F_{\mathcal Q}.
\]
The quotient \(p\) induces a continuous functor
\[
\widehat p:\widehat F\to P.
\]

Since \(p\) satisfies every equation valid in \(\mathcal Q\), any two
profinite arrows of \(\widehat F\) with the same image in
\[
\widehat F_{\mathcal Q}
\]
have the same image under
\[
\widehat p.
\]
Equivalently,
\[
\ker(\pi_{\mathcal Q})\leq\ker(\widehat p).
\]
Therefore
\[
\widehat p
\]
factors through a continuous functor
\[
\bar p:\widehat F_{\mathcal Q}\to P.
\]

By Proposition~\ref{prop:reiterman-finite-maps-factor-subdirect}, the functor
\[
\bar p
\]
factors through a finite subdirect product
\[
r:F\twoheadrightarrow R
\]
of quotients in \(\mathcal Q\). Since \(\mathcal Q\) is closed under finite
subdirect products,
\[
r\in\mathcal Q.
\]
Since
\[
p
\]
is a further finite quotient of
\[
r,
\]
closure under further finite quotients gives
\[
p\in\mathcal Q.
\]
Thus
\[
\mathcal Q
=
\operatorname{Mod}_{\mathcal R^{\mathrm{fin}}}(\operatorname{Eq}(\mathcal Q)).
\]
\end{proof}

\begin{remark}[Why this is separate from intersection compactness]
\label{rem:reiterman-why-finite-separate}
The proof of Theorem~\ref{thm:reiterman-finite-profinite} does not assert
that finite quotients are intersection-compact. It uses a different
compactness principle: every continuous map from a profinite inverse limit to
a finite discrete object factors through finitely many quotient coordinates.
This is the compactness principle underlying classical Reiterman theory.
\end{remark}

\section{Global finite/profinite Reiterman theorem}
\label{sec:reiterman-global-finite-profinite}

We now combine the finite/profinite theorem with Reiterman-admissible
interface translations.

Assume
\[
\mathcal E=\mathbf{Set}.
\]
Let
\[
\mathfrak T
\]
be a category of interfaces and Reiterman-admissible translations. For each
interface
\[
\Xi
\]
let
\[
\mathcal R_\Xi^{\mathrm{fin}}
\]
be a finite quotient collection for
\[
F_\Xi.
\]
Assume the object set of each \(F_\Xi\) is finite, so that the
identity-on-objects quotients in
\[
\mathcal R_\Xi^{\mathrm{fin}}
\]
are finite internal categories precisely when their arrow sets are finite.

Assume that Reiterman-admissible translations preserve finite quotient
collections:
\[
q\in\mathcal R_\Xi^{\mathrm{fin}}
\quad\Longrightarrow\quad
\Gamma_{\mathrm{alg}}^\ast q\in\mathcal R_{\Xi'}^{\mathrm{fin}}
\]
for every
\[
\Gamma:\Xi'\to\Xi.
\]

\begin{definition}[Global finite quotient variety]
\label{def:reiterman-global-finite-quotient-variety}
A \textbf{global finite algebraic quotient variety} is a family
\[
\mathcal Q=(\mathcal Q_\Xi)_\Xi,
\qquad
\mathcal Q_\Xi\subseteq\mathcal R_\Xi^{\mathrm{fin}},
\]
such that:
\begin{enumerate}
\item each \(\mathcal Q_\Xi\) is closed under isomorphism of quotients under
\(F_\Xi\);
\item each \(\mathcal Q_\Xi\) is closed under further finite quotients;
\item each \(\mathcal Q_\Xi\) is closed under finite subdirect products;
\item the family is stable under Reiterman-admissible translations:
\[
q\in\mathcal Q_\Xi
\quad\Longrightarrow\quad
\Gamma_{\mathrm{alg}}^\ast q\in\mathcal Q_{\Xi'}.
\]
\end{enumerate}
\end{definition}

\begin{definition}[Global profinite equational theory]
\label{def:reiterman-global-profinite-equational-theory}
A \textbf{global profinite equational theory}
\[
\Sigma=(\Sigma_\Xi)_\Xi
\]
assigns to every interface
\[
\Xi
\]
a class
\[
\Sigma_\Xi
\]
of profinite completed input equations over
\[
\widehat F_\Xi^{\,\mathrm{fin}}
\]
and is model-stable under Reiterman-admissible translations:
\[
q\models\Sigma_\Xi
\quad\Longrightarrow\quad
\Gamma_{\mathrm{alg}}^\ast q\models\Sigma_{\Xi'}
\]
for every
\[
\Gamma:\Xi'\to\Xi.
\]
\end{definition}

\begin{theorem}[Global finite/profinite Reiterman theorem]
\label{thm:reiterman-global-finite-profinite}
Let
\[
(\mathcal R_\Xi^{\mathrm{fin}})_\Xi
\]
be a translation-stable family of finite quotient collections. For a family
\[
\mathcal Q=(\mathcal Q_\Xi)_\Xi,
\qquad
\mathcal Q_\Xi\subseteq\mathcal R_\Xi^{\mathrm{fin}},
\]
the following are equivalent:
\begin{enumerate}
\item \(\mathcal Q\) is a global finite algebraic quotient variety.
\item There exists a global profinite equational theory
\[
\Sigma=(\Sigma_\Xi)_\Xi
\]
such that
\[
\mathcal Q_\Xi
=
\operatorname{Mod}_{\mathcal R_\Xi^{\mathrm{fin}}}(\Sigma_\Xi)
\]
for every interface \(\Xi\).
\end{enumerate}
Moreover, one may take
\[
\Sigma_\Xi
\]
to be the class of all profinite equations valid in
\[
\mathcal Q_\Xi.
\]
\end{theorem}

\begin{proof}
The local equivalence at each interface is
Theorem~\ref{thm:reiterman-finite-profinite}.

If
\[
\mathcal Q
=
\operatorname{Mod}_{\mathcal R^{\mathrm{fin}}}(\Sigma)
\]
for a global profinite equational theory, then the local closure properties
follow objectwise from the finite/profinite theorem. Translation stability
follows from the model-stability condition in
Definition~\ref{def:reiterman-global-profinite-equational-theory}.

Conversely, suppose
\[
\mathcal Q
\]
is a global finite algebraic quotient variety. Objectwise, define
\[
\Sigma_\Xi
\]
to be the class of all profinite completed equations valid in
\[
\mathcal Q_\Xi.
\]
The local finite/profinite theorem gives
\[
\mathcal Q_\Xi
=
\operatorname{Mod}_{\mathcal R_\Xi^{\mathrm{fin}}}(\Sigma_\Xi).
\]
To see model-stability, let
\[
q\models\Sigma_\Xi.
\]
Then by the local theorem,
\[
q\in\mathcal Q_\Xi.
\]
For a Reiterman-admissible translation
\[
\Gamma:\Xi'\to\Xi,
\]
global translation stability gives
\[
\Gamma_{\mathrm{alg}}^\ast q\in\mathcal Q_{\Xi'}.
\]
Again by the local theorem,
\[
\Gamma_{\mathrm{alg}}^\ast q\models\Sigma_{\Xi'}.
\]
Therefore \(\Sigma\) is a global profinite equational theory and its model
class is exactly \(\mathcal Q\).
\end{proof}

\section{The one-object Set specialization}
\label{sec:reiterman-one-object-set}

Let
\[
\mathcal E=\mathbf{Set},
\qquad
\mathbb A=M,
\qquad
\mathbb B=N,
\]
where \(M\) and \(N\) are monoids. We retain the convention
\[
mn=m\circ n.
\]
Thus \(mn\) means first \(n\), then \(m\) as categorical composition, while
right-action execution by \(mn\) is computed as first \(m\), then \(n\).

Let
\[
D\subseteq\mathsf{Beh}_{M,N}
\]
be residually closed. Then
\[
\operatorname{TransOut}_{N}(D)
\]
is the one-object category whose monoid consists of pairs
\[
(o,\tau),
\]
with
\[
\tau:D\to D,
\qquad
o:D\to N.
\]
Multiplication is
\[
(o,\tau)(o',\tau')
=
\left(
\lambda\mapsto o(\lambda)o'(\tau\lambda),
\ \tau'\circ\tau
\right),
\]
and the unit is
\[
(\lambda\mapsto 1_N,\operatorname{id}_D).
\]

The canonical representation is
\[
M\to \operatorname{TransOut}_N(D),
\qquad
m\longmapsto(o_m,\tau_m),
\]
where
\[
\tau_m(\lambda)=\partial_m\lambda
\]
and
\[
o_m(\lambda)=\omega_D(m,\lambda).
\]
If
\[
\lambda=\partial_r\lambda_0
\]
with
\[
\lambda_0\in S,
\]
then
\[
o_m(\partial_r\lambda_0)=\lambda_0(r,m).
\]

Categorically, this representation is the one-object instance of
\[
\mathbb A^{\mathrm{op}}
\to
\operatorname{TransOut}_{\mathbb B}(D).
\]
With the operational multiplication convention
\[
mn=m\circ n,
\]
it is written as a monoid homomorphism
\[
M\to\operatorname{TransOut}_N(D).
\]

\begin{proposition}[One-object syntactic algebraic recognizer]
\label{prop:reiterman-one-object-synalg}
For
\[
S\subseteq\mathsf{Beh}_{M,N},
\]
the general syntactic algebraic recognizer specializes to
\[
\operatorname{SynAlg}_{M,N}(S)
=
\operatorname{im}
\left(
M\to
\operatorname{TransOut}_N(\mathsf D(S))
\right).
\]
\end{proposition}

\begin{proof}
In the one-object case, the internal category
\[
\operatorname{TransOut}_{\mathbb B}(\mathsf D(S))
\]
has one object, and its monoid of arrows is precisely
\[
\operatorname{TransOut}_N(\mathsf D(S)).
\]
The functor
\[
\rho_{\mathsf D(S)}:
\mathbb A^{\mathrm{op}}\to
\operatorname{TransOut}_{\mathbb B}(\mathsf D(S))
\]
is exactly the monoid homomorphism
\[
M\to
\operatorname{TransOut}_N(\mathsf D(S))
\]
described above. Taking its image gives the displayed formula.
\end{proof}

\section{The free-input specialization}
\label{sec:reiterman-free-input-specialization}

Suppose now that
\[
M=I^\ast
\]
is the free monoid on a set \(I\). As in the preceding chapter, the final
behaviour object has the generator-level normal form
\[
\mathsf{Beh}_{I^\ast,N}
\cong
N^{I^+},
\]
and for
\[
S\subseteq N^{I^+},
\]
the residual closure is
\[
\mathsf D(S)
=
\{\partial_w\beta\mid w\in I^\ast,\ \beta\in S\}.
\]

The syntactic algebraic recognizer is
\[
\operatorname{SynAlg}_{I^\ast,N}(S)
=
\operatorname{im}
\left(
I^\ast
\to
\operatorname{TransOut}_N(\mathsf D(S))
\right).
\]

\begin{definition}[Finite syntactic specification]
\label{def:reiterman-finite-syntactic-specification}
A specification
\[
S\subseteq N^{I^+}
\]
is \textbf{finite syntactic} if its syntactic input quotient
\[
I^\ast
\twoheadrightarrow
\operatorname{SynAlg}_{I^\ast,N}(S)
\]
is finite.
\end{definition}

\begin{remark}[Finite state as a special case]
\label{rem:reiterman-finite-state-special-case}
When \(I\), \(N\), and the residual state set \(\mathsf D(S)\) are finite,
the syntactic algebraic recognizer is finite, giving the usual finite
quotient situation of classical Set-based recognition. In the abstract
intersection-compact theorem, finiteness is not identified with
intersection-compactness. In the finite/profinite theorem, finiteness is
handled instead through profinite compactness and finite determinacy.
\end{remark}

\section{The Eilenberg--Reiterman triangle}
\label{sec:reiterman-triangle}

The Eilenberg--Reiterman picture has three levels:
\[
\text{behaviour specifications}
\quad\longleftrightarrow\quad
\text{residual state recognizers}
\quad\longrightarrow\quad
\text{effective input quotients}.
\]

The first correspondence is the state Eilenberg correspondence:
\[
\text{stable classes of closed \(\mathcal K\)-state recognizers}
\quad\cong\quad
\text{\(\mathcal K\)-recognizable behaviour varieties}.
\]
This was established in the preceding chapter.

The second passage sends a closed state recognizer \(D\) to the effective
input quotient
\[
\mathbb A^{\mathrm{op}}
\twoheadrightarrow
\operatorname{Alg}_{\mathbb A,\mathbb B}(D)
=
\operatorname{im}
\left(
\mathbb A^{\mathrm{op}}
\to
\operatorname{TransOut}_{\mathbb B}(D)
\right).
\]
For a behaviour specification \(S\), this gives the syntactic algebraic
recognizer
\[
\operatorname{SynAlg}_{\mathbb A,\mathbb B}(S)
=
\operatorname{Alg}_{\mathbb A,\mathbb B}(\mathsf D(S)).
\]

Globally, Reiterman-admissible interface translations satisfy
\[
\Gamma_{\mathrm{beh}}^\ast\mathsf D_\Xi
=
\mathsf D_{\Xi'}\Gamma_{\mathrm{beh}}^\ast
\]
and
\[
\Gamma_{\mathrm{alg}}^\ast
\operatorname{Alg}_\Xi(D)
\cong
\operatorname{Alg}_{\Xi'}(\Gamma_{\mathrm{st}}^\ast D).
\]
Consequently,
\[
\Gamma_{\mathrm{alg}}^\ast
\operatorname{SynAlg}_\Xi(S)
\cong
\operatorname{SynAlg}_{\Xi'}(\Gamma_{\mathrm{beh}}^\ast S).
\]
This is the bridge from global behaviour varieties to global algebraic
quotient classes.

The Reiterman correspondence belongs to the effective-input-quotient level.
There are two forms. The abstract form uses \(\kappa\)-intersection-compact
quotients. The finite Set-based form uses profinite completed input
categories and finite determinacy.

Thus the state recognizer
\[
\mathsf D(S)
\]
is the minimal separated residual realization of \(S\), while
\[
\operatorname{SynAlg}_{\mathbb A,\mathbb B}(S)
\]
is the algebra of all input-induced transition-output operations on that
minimal realization.

\section{Main consolidation}
\label{sec:reiterman-main-consolidation}

\begin{theorem}[Reiterman theory for effective input quotients of residual Mealy recognizers]
\label{thm:reiterman-main-consolidation}
Let
\[
\mathcal E
\]
be a Grothendieck topos, and let
\[
\mathbb A,\mathbb B
\]
be internal categories in \(\mathcal E\). Let
\[
\mathsf{Beh}_{\mathbb A,\mathbb B}
\]
be the final behaviour system for Mealy morphisms
\[
\mathbb A\rightsquigarrow\mathbb B.
\]
Then:

\begin{enumerate}
\item For every closed behaviour system
\[
D\hookrightarrow\mathsf{Beh}_{\mathbb A,\mathbb B},
\]
there is an internal transition-output endomorphism category
\[
\operatorname{TransOut}_{\mathbb B}(D)
\]
with object object \(A_0\).

\item The Mealy structure of \(D\) induces an internal functor
\[
\rho_D:
\mathbb A^{\mathrm{op}}
\to
\operatorname{TransOut}_{\mathbb B}(D).
\]

\item The algebraic recognizer of \(D\) is the image
\[
\operatorname{Alg}_{\mathbb A,\mathbb B}(D)
=
\operatorname{im}(\rho_D).
\]

\item For a behaviour specification \(S\), the syntactic algebraic recognizer
is
\[
\operatorname{SynAlg}_{\mathbb A,\mathbb B}(S)
=
\operatorname{Alg}_{\mathbb A,\mathbb B}(\mathsf D(S)).
\]

\item Equivalently,
\[
\operatorname{SynAlg}_{\mathbb A,\mathbb B}(S)
\cong
\mathbb A^{\mathrm{op}}/{\equiv_{\mathsf D(S)}},
\]
where two input arrows are equivalent precisely when they induce the same
transition and output operation on all residual states of \(\mathsf D(S)\).

\item If a Mealy system \(X\) recognizes \(S\), then
\[
\operatorname{SynAlg}_{\mathbb A,\mathbb B}(S)
\]
is a regular quotient of
\[
\operatorname{Alg}_{\mathbb A,\mathbb B}(\operatorname{Real}(X)).
\]

\item Effective input quotients of
\[
F=\mathbb A^{\mathrm{op}}
\]
are equivalently effective internal category congruences on \(F\).

\item Subdirect products of effective input quotients correspond to
intersections of their kernel congruences.

\item Classes of effective input quotients closed under further quotients and
arbitrary small subdirect products are exactly the classes definable by
generalized completed input equations in the completed input category
determined by those quotients.

\item For any regular cardinal \(\kappa\), classes of
\(\kappa\)-intersection-compact effective input quotients closed under
further compact quotients and \(<\kappa\)-subdirect products are exactly the
classes definable by generalized completed input equations in the
corresponding \(\kappa\)-compact completed input category.

\item Under a Reiterman-admissible interface translation
\[
\Gamma:\Xi'\to\Xi,
\]
syntactic algebraic recognizers satisfy
\[
\Gamma_{\mathrm{alg}}^\ast
\operatorname{SynAlg}_{\Xi}(S)
\cong
\operatorname{SynAlg}_{\Xi'}(\Gamma_{\mathrm{beh}}^\ast S).
\]

\item Global algebraic quotient classes closed locally under the appropriate
quotient and subdirect-product operations and globally under
Reiterman-admissible translations are exactly the classes definable by global
completed equational theories.

\item In the Set-based finite case, finite quotient varieties are described
by profinite completed equations. The proof uses compactness of profinite
input completions and finite determinacy of continuous maps into finite
discrete quotients, not intersection compactness of finite kernels.

\item In the one-object Set case
\[
\mathbb A=M,
\qquad
\mathbb B=N,
\]
the syntactic algebraic recognizer is
\[
\operatorname{SynAlg}_{M,N}(S)
=
\operatorname{im}
\left(
M\to
\operatorname{TransOut}_N(\mathsf D(S))
\right).
\]

\item If
\[
M=I^\ast,
\]
then
\[
\operatorname{SynAlg}_{I^\ast,N}(S)
=
\operatorname{im}
\left(
I^\ast
\to
\operatorname{TransOut}_N(\mathsf D(S))
\right).
\]
\end{enumerate}
\end{theorem}

\begin{proof}
Parts (1) and (2) are
Definition~\ref{def:reiterman-transout-category},
Definition~\ref{def:reiterman-rho-D},
Proposition~\ref{prop:reiterman-transout-is-internal-category}, and
Proposition~\ref{prop:reiterman-rho-D-functor}.
Parts (3)--(5) are
Definitions~\ref{def:reiterman-alg-D},
\ref{def:reiterman-synalg}, and
\ref{def:reiterman-syntactic-input-congruence}, together with
Proposition~\ref{prop:reiterman-syntactic-input-congruence}.
Part (6) is
Theorem~\ref{thm:reiterman-algebraic-syntactic-divisibility}.
Parts (7) and (8) are
Proposition~\ref{prop:reiterman-congruences-effective-quotients} and
Proposition~\ref{prop:reiterman-kernel-subdirect-product}.
Part (9) is
Theorem~\ref{thm:reiterman-complete-equational-theorem}.
Part (10) is
Theorem~\ref{thm:reiterman-intersection-compact}.
Part (11) is
Proposition~\ref{prop:reiterman-synalg-commutes-with-translations}.
Part (12) is
Theorem~\ref{thm:reiterman-global-intersection-compact}.
Part (13) is
Theorem~\ref{thm:reiterman-finite-profinite} and
Theorem~\ref{thm:reiterman-global-finite-profinite}.
Part (14) is
Proposition~\ref{prop:reiterman-one-object-synalg}.
Part (15) is the free-input specialization of the same construction.
\end{proof}